

% ------------------------------------------------------------------------------
% Fundamentação Teórica
% ------------------------------------------------------------------------------

\chapter{Fundamentação Teórica}
\label{chap_fundamentacao_teorica}

A fundamentação teórica apresenta os conceitos físicos e matemáticos necessários para compreender o fenômeno da queda livre e os modelos utilizados na simulação computacional. São discutidos: a dinâmica da queda livre sem resistência do ar, a distinção entre partícula e corpo rígido, e os modelos físicos de resistência do ar.

\section{Queda Livre Sem Resistência do Ar}

A queda livre é o movimento de um corpo sujeito exclusivamente à força da gravidade. Nas proximidades da superfície da Terra, considera-se a aceleração gravitacional constante, denotada por $g$.

A dinâmica do movimento é descrita pelas equações diferenciais:
\begin{equation}
	a(t) = -g
	\label{eq:aceleracao}
\end{equation}
\begin{equation}
	v(t) = v_0 - g t
	\label{eq:velocidade}
\end{equation}
\begin{equation}
	y(t) = H + v_0 t - \frac{1}{2} g t^2
	\label{eq:posicao}
\end{equation}
onde $H$ é a altura inicial e $v_0$ a velocidade inicial.

\subsection*{Aproximação Discreta (Modelo Numérico)}

Para implementar numericamente o movimento, utiliza-se a discretização temporal. O tempo é dividido em intervalos de tamanho $\Delta t$, e as equações diferenciais são aproximadas por diferenças finitas:
\begin{equation}
	t_{i+1} = t_i + \Delta t
	\label{eq:tempo_discreto}
\end{equation}
\begin{equation}
	y_{i+1} = y_i + v_i \Delta t
	\label{eq:posicao_discreta}
\end{equation}
\begin{equation}
	v_{i+1} = v_i - g \Delta t
	\label{eq:velocidade_discreta}
\end{equation}
Essas fórmulas (Eqs.~\ref{eq:tempo_discreto}--\ref{eq:velocidade_discreta}) permitem simular o movimento passo a passo até que a posição atinja o solo ($y_i \le 0$).

\section{Partícula versus Corpo Rígido}

No modelo básico, o corpo é tratado como uma \textbf{partícula}, ou seja, um ponto sem dimensões, sem forma e sem rotação. Essa simplificação é válida quando não há resistência do ar e apenas a força peso atua, sendo o interesse no movimento do centro de massa.

Um \textbf{corpo rígido}, por outro lado, possui forma geométrica, área frontal, distribuição de massa, possibilidade de rotação e interação aerodinâmica com o ar. Sem resistência do ar, o centro de massa de um corpo rígido obedece exatamente às mesmas equações da partícula (Eqs.~\ref{eq:aceleracao}--\ref{eq:posicao}). Entretanto, ao considerar o ar, a forma e a orientação do corpo influenciam o movimento, exigindo modelos mais completos.

\section{Resistência do Ar}

Ao abandonar a hipótese de partícula ideal e considerar um corpo real, surge a necessidade de incluir a \textbf{força de arrasto}, que se opõe ao movimento. A resistência do ar depende da velocidade, da forma do corpo e das propriedades do fluido.

\subsection{Modelo Linear de Arrasto}

Adequado para baixas velocidades ou fluidos muito viscosos:
\begin{equation}
	F_d = -k v
	\label{eq:arrasto_linear}
\end{equation}
onde $k$ é o coeficiente de arrasto linear. A equação do movimento torna-se:
\begin{equation}
	m \frac{dv}{dt} = -mg - k v
	\label{eq:movimento_linear}
\end{equation}

\subsection{Modelo Quadrático de Arrasto}

Modelo fisicamente mais realista para objetos caindo no ar:
\begin{equation}
	F_d = -c v |v|
	\label{eq:arrasto_quadratico}
\end{equation}
ou, na forma aerodinâmica completa:
\begin{equation}
	F_d = \frac{1}{2} \rho C_d A v^2
	\label{eq:arrasto_aerodinamico}
\end{equation}
onde $\rho$ é a densidade do ar, $C_d$ o coeficiente aerodinâmico e $A$ a área frontal do corpo. Esse modelo prevê a existência de uma \textbf{velocidade terminal}, quando o arrasto equilibra o peso.

\section{Considerações sobre Acurácia Numérica}

A simulação computacional utiliza aproximações discretas, o que introduz erro numérico. A precisão do método depende principalmente do tamanho do passo de tempo $\Delta t$ e da forma como as derivadas são aproximadas. Reduzir $\Delta t$ melhora a aproximação das equações contínuas e diminui o erro acumulado. Métodos mais avançados podem utilizar valores intermediários de velocidade para aumentar a precisão, mas o princípio fundamental permanece: quanto menor o passo, maior a acurácia.


\section{Queda Livre Sem Resistência do Ar}

A queda livre clássica considera que a única força atuante é a gravidade. Nesse caso, o movimento é uniformemente acelerado, com aceleração constante igual a $g$. Não há qualquer termo de resistência do ar. Esse modelo é válido para corpos pequenos, densos ou em situações onde o efeito do ar é desprezível.

A equação diferencial fundamental desse movimento é:
\begin{equation}
	m \frac{d^2y}{dt^2} = -mg
	\label{eq:movimento}
\end{equation}
onde $y$ é a posição vertical, $m$ a massa do corpo e $g$ a aceleração da gravidade. A solução analítica para a posição em função do tempo, considerando altura inicial $H$ e velocidade inicial $v_0$, é:
\begin{equation}
	y(t) = H + v_0 t - \frac{1}{2}gt^2
	\label{eq:posicao}
\end{equation}
e para a velocidade:
\begin{equation}
	v(t) = v_0 - g t
	\label{eq:velocidade}
\end{equation}
e para a aceleração:
\begin{equation}
	a(t) = -g
	\label{eq:aceleracao}
\end{equation}

Essas expressões (Eqs.~\ref{eq:posicao}, \ref{eq:velocidade} e \ref{eq:aceleracao}) descrevem completamente o movimento de queda livre sem resistência do ar.

\subsection*{Aproximação Discreta (Modelo Numérico)}

Para simulação computacional, as equações diferenciais são aproximadas por diferenças finitas, discretizando o tempo em intervalos de tamanho $\Delta t$:
\begin{equation}
	t_{i+1} = t_i + \Delta t
	\label{eq:tempo_discreto}
\end{equation}
\begin{equation}
	y_{i+1} = y_i + v_i \Delta t
	\label{eq:posicao_discreta}
\end{equation}
\begin{equation}
	v_{i+1} = v_i - g \Delta t
	\label{eq:velocidade_discreta}
\end{equation}
Essas fórmulas (Eqs.~\ref{eq:tempo_discreto}--\ref{eq:velocidade_discreta}) permitem calcular iterativamente a trajetória do corpo até atingir o solo.




\section{Queda Livre com Resistência do Ar}

Quando o corpo possui área significativa, formato não esférico ou velocidade elevada, a resistência do ar não pode ser desprezada. Existem dois modelos principais para o arrasto:

\subsection{Resistência Linear}
Ocorre para baixas velocidades ou objetos pequenos, onde a força de resistência é proporcional à velocidade:
\begin{equation}
	m \frac{dv}{dt} = -mg - c v
	\label{eq:linear}
\end{equation}
onde $c$ é o coeficiente de arrasto linear. A solução dessa equação mostra que a velocidade terminal é atingida quando a força de resistência iguala o peso.

\subsection{Resistência Quadrática}
Para velocidades mais altas ou corpos maiores, a força de resistência é proporcional ao quadrado da velocidade:
\begin{equation}
	m \frac{dv}{dt} = -mg - k v^2
	\label{eq:quadratica}
\end{equation}
onde $k$ é o coeficiente de arrasto quadrático, dependente da área, formato e densidade do ar. Esse modelo é mais realista para objetos em queda livre no ar. A solução dessa equação também leva a uma velocidade terminal, diferente do caso linear.



\section{Influência das Características do Corpo: Corpo Rígido em Queda Livre}

Ao considerar um corpo rígido, a modelagem se torna mais realista e abrangente. A resistência do ar passa a ser relevante para corpos com área significativa ou formas não esféricas \cite[p.~49]{halliday2011}. Além disso, a rotação pode alterar a área de contato com o ar, modificando o arrasto e tornando o movimento mais complexo. Em situações práticas, a velocidade terminal depende das características do corpo, como massa, área, formato e coeficiente de arrasto, exigindo abordagens matemáticas e computacionais mais sofisticadas. Assim, a escolha entre resistência linear (Eq.~\ref{eq:linear}) ou quadrática (Eq.~\ref{eq:quadratica}) depende do regime de movimento e das propriedades do corpo.



\section{Resumo dos Efeitos da Resistência do Ar}

Batchelor afirma que “a força de resistência pode ser proporcional à velocidade ou ao quadrado da velocidade, dependendo do regime de escoamento” \cite[p.~246]{batchelor2000}. A inclusão desses termos (Eqs.~\ref{eq:linear} e \ref{eq:quadratica}) altera significativamente o comportamento do movimento, reduzindo a velocidade terminal e aumentando o tempo de queda. A análise desses efeitos é fundamental para aplicações práticas e experimentais.




Por fim, vale destacar que, quando a solução analítica não é viável devido à complexidade do modelo ou à presença de forças adicionais, recorre-se a métodos numéricos. O método de Euler, por exemplo, é amplamente utilizado por sua simplicidade e eficiência na resolução de equações diferenciais ordinárias \cite[p.~7]{butcher2016}. O uso dessas técnicas permite simular cenários mais realistas, tornando possível a análise computacional de sistemas físicos complexos, como a queda livre com resistência do ar ou corpos de geometria não trivial.
